% letter-formal.tex — a formal letter using the standard letter class, with a
% second letter in the same file to show how the class handles more than one.
% Compile with: pdflatex letter-formal.tex
%
% Both correspondents, the addresses, and the reference numbers are invented.

\documentclass[11pt,a4paper]{letter}

\usepackage[utf8]{inputenc}
\usepackage[T1]{fontenc}
\usepackage{lmodern}
\usepackage{geometry}

\geometry{margin=2.5cm}

\signature{Marguerite Oyelaran\\Head of Collections}
\address{
  Fenwick District Libraries \\
  118 Marlow Street \\
  Fenwick FW2 4RG
}

\begin{document}

\begin{letter}{
  Mr. D.\ Ashworth \\
  Ashworth \& Pell Bindery \\
  7 Cooper's Yard \\
  Fenwick FW1 9BT
}

\opening{Dear Mr. Ashworth,}

\textbf{Re: Conservation of the Marlow Street bequest --- reference FDL/27/0412}

Thank you for your estimate of 3 August and for the time your colleague spent
with our collections officer last week.

The Board has approved the work in full, at the quoted figure of £14,600
excluding VAT, and I am writing to confirm the terms we discussed so that both
of us are working from the same understanding.

The bequest comprises 212 volumes. Of these, 47 require full rebinding, 96
require repair to the spine or boards only, and the remaining 69 require
cleaning and boxing. Your estimate treats these as three separate schedules,
which suits us; we would like to be invoiced against each schedule as it
completes rather than on a single invoice at the end.

Three points were left open at our meeting, and I should set out our position on
each:

\begin{enumerate}
  \item \textbf{Original materials.} Where a board or an endpaper can be
        retained, it should be, even at the cost of a less uniform finish
        across the set. Uniformity is not what this collection is for.
  \item \textbf{Documentation.} We ask for a photographic record before and
        after treatment of every volume in the first two schedules, supplied as
        files rather than prints.
  \item \textbf{Transport.} Volumes will move in batches of no more than
        twenty, and will remain insured under our policy while in your
        possession. Please confirm that this is acceptable to your own insurer.
\end{enumerate}

Subject to your confirmation of the third point, we would like work to begin on
14 September. Our collections officer, Mr.\ Tobias Renn, will be your contact
for scheduling and will arrange the first collection.

Please countersign the enclosed schedule of works and return one copy.

\closing{Yours sincerely,}

\encl{Schedule of works (two copies)}
\cc{T.\ Renn, Collections Officer \\ Finance Office, Fenwick District Council}
\ps{P.S. The reading room will be closed for its own refurbishment from 2 to 16
September. Deliveries during that fortnight should go to the Cooper Street
annexe instead.}

\end{letter}

\begin{letter}{
  Ms. P.\ Lindqvist \\
  Fenwick District Council \\
  Finance Office \\
  Fenwick FW1 1AA
}

\opening{Dear Ms. Lindqvist,}

\textbf{Re: Commitment against budget line 4420 --- reference FDL/27/0413}

Further to the Board's decision of 11 August, I am writing to commit £14,600
excluding VAT against budget line 4420 for the conservation of the Marlow Street
bequest.

The work will be invoiced in three parts across the current and following
financial quarters, as set out in the attached schedule. No part of the
commitment falls outside the current financial year.

The bequest was accepted on the condition that it be made available to readers
within two years of receipt. That deadline falls in June of next year, and the
conservation work is the only outstanding obstacle to meeting it.

I am happy to answer any questions the Finance Office may have.

\closing{Yours sincerely,}

\encl{Payment schedule}

\end{letter}

\end{document}
